\section{The Relativistic Orbit Coefficient}

Boundary radiation was already forced before this chapter began. The
interior mixed term had been forbidden; the terminal residue had to be read
at the boundary. Chapter 09 then carried alpha to that boundary as a finite
receipt. The new pressure is that the boundary is no longer an ordinary
edge. It is being pushed toward the lightlike face, where the cost of
exporting an interior refinement no longer behaves like the cost of an
ordinary timelike comparison.

The orbit coefficient names the cost of preserving the invariant as the
surface reading is pressed toward that face. It is not inserted as a
coefficient by taste. A finite orbit can be bracketed by three readings:
\[
  r_- < r_0 < r_+,
  \qquad
  v_-^2 > 1,\quad v_0^2 = 1,\quad v_+^2 < 1.
\]
The lower side is beyond the ordinary timelike reading. The middle face is
lightlike. The upper side is timelike and therefore still carries a finite
approach reading. The grammar keeps the bracket because each face says what
kind of comparison is still admissible.

The lightlike face is not an interior state waiting to be inspected. It is
the boundary at which the export grammar changes. A timelike reading can
carry a finite coefficient for the approach. A lightlike reading marks the
surface where the coefficient has become the price of staying on the
boundary rather than a property of an interior orbit. To read it as an
ordinary orbit would be to smuggle an interior path through the horizon.

The radiation residue is therefore pushed, not crossed. The boundary trace
can be refined toward the null surface. The residue can be oriented by the
coefficient that tells how expensive the approach has become. But the
grammar forbids the reader from treating the lightlike face as if the
interior had been opened. The approach names a direction of pressure, not a
completed possession of what lies beyond the surface.

This is the first quantum-gravity move of the chapter. Gravity appears here
as the spin-residue cost of pushing a surface reading toward a null horizon.
Quantum appears as the finite refusal to replace that residue by a smooth
interior story. The coefficient carries both obligations. It reads how the
boundary radiation has been forced toward the lightlike face, and it
forbids the face from becoming an unbounded interior measurement.

A horizon example clarifies the rule. An exterior ledger may continue to
receive surface refinements while the interior continues to refine on its
own side. Near the horizon, the exchange rate between those ledgers
collapses. The exterior can still record a boundary correction. It cannot
merge every interior update into its own finite history. The orbit
coefficient is the cost-bearing reading of that collapse as a surface
approach.

The coefficient also protects alpha. Alpha enters this chapter as a receipt
with path cost attached. If the surface reading approaches a lightlike
horizon, the receipt cannot be detached from that approach and treated as a
bare number. Its value, inverse reading, and correction residue are carried
on the surface together with the cost of the approach. The horizon may mask
the interior, but it may not erase the receipt history that makes the
surface valuation lawful.

Thus the relativistic orbit coefficient does not add a new substance to the
ledger. It changes the status of the boundary. The radiation residue is
pushed toward a null surface; the orbit coefficient carries the cost of that
push; the grammar permits the approach and forbids a completed crossing. The
next section turns that approach into a certificate.
